\documentclass{article}

\textwidth=6.5in
\textheight=8.5in
\voffset=-54pt
\oddsidemargin=0in
\evensidemargin=0in
\setlength{\parskip}{8pt}
\setlength{\parindent}{0pt}

\usepackage{workbook}

\usepackage[sols]{handout}

\begin{document}

\begin{center}
  \large\textsc{Problem Set 33 - More Density Problems}
  
      \pspace{12pt}
  \begin{problemonly}\large Name: \rule{5cm}{0.15mm}\\ \end{problemonly}
\end{center}

\begin{grading}
  
    \begin{enumerate}[nosep]
    \item 15 points 
    \item 20 points 
    \item 10 points
    \item 2 points
    \item 5 point
    \end{enumerate}

    Total: 52 points
  
\end{grading}

\begin{enumerate}
\item

For help with this problem, read pp. 827-838 in Gottlieb.

  \begin{enumerate}


  \item
\begin{grading}
    15 points. Part (a): 2/3/1/1/3. Part (b): 2/2/1
\end{grading}
  
  Suppose that the density of organisms in a circular petri dish
    varies with distance from the center point of the dish.  The density
    at a distance $x$ centimeters from the center is given by $f(x)$
    organisms per square centimeter. The petri dish is 18 centimeters in
    diameter.
    \begin{enumerate}
    \item
      \begin{problem}[\stretch{3}]
        Show in a sketch how you will slice the petri dish to solve this
        problem.
      \end{problem}

      \begin{solution}
        Since the density varies with distance to the center, we should
        slice the dish in concentric rings.  Here's a picture, which also
        shows the meaning of $x$; since $x$ measures distance to the
        center, $x = 0$ at the center and $x = 9$ at the edge of the petri
        dish:
        \begin{center}
          \begin{tikzpicture}
            \begin{axis}[x=0.007in, y=0.007in, xmin=0, xmax=115,
              ymin=-100, ymax=100, axis on top=true, clip=false, hide y
              axis, xtick={0.1, 100}, xticklabels={0, 9},
              xlabel=$x$, tick label style={fill=none}]              
              %\draw[fill=cyan, draw=none] (0, 0) circle(75); %
              %\draw[fill=white, draw=none] (0, 0) circle (62.5); %
              \pgfplotsforeachungrouped \x in {12.5, 25, ..., 100.1}
              {
                \edef\temp{
                  \noexpand\draw (0, 0) circle(\x);
                }
                \temp
              }              
            \end{axis}
          \end{tikzpicture}
        \end{center}
        Therefore, we're slicing the interval $[0, 9]$. 
      \end{solution}

    \item
      \begin{problem}[\vfill]
        Approximate the number of organisms in the $k$-th slice.
      \end{problem}

      \begin{solution}
        If we label $x_0 = 0, \ldots, x_n = 9$ as usual, here's a picture
        of the $k$-th slice:
        \begin{center}
          \begin{tikzpicture}
            \begin{axis}[x=0.007in, y=0.007in, xmin=0, xmax=115,
              ymin=-100, ymax=100, axis on top=true, clip=false, hide y
              axis, xtick={0.1, 12.5, 25, ..., 100},
              xticklabels={$x_0$,,,,,,$x_k$,,$x_n$}, 
              xlabel=$x$, tick label style={fill=none}]              
              \draw[fill=cyan, draw=none] (0, 0) circle(75); %
              \draw[fill=white, draw=none] (0, 0) circle (62.5); %
              \pgfplotsforeachungrouped \x in {12.5, 25, ..., 100.1}
              {
                \edef\temp{
                  \noexpand\draw (0, 0) circle(\x);
                }
                \temp
              }              
            \end{axis}
          \end{tikzpicture}
        \end{center}
        The slice is shaped like a ring with outer radius $x_k$ and
        thickness $\Delta x$, so its area is approximately $2 \pi x_k
        \Delta x$.  The density of organisms in this slice is
        approximately $f(x_k)$, so the number of organisms in this slice
        is approximately $(\text{density}) \cdot (\text{area}) \approx
        \boxed{f(x_k) \cdot 2 \pi x_k \Delta x}$.
      \end{solution}

    \item
      \begin{problem}[\vfill]
        Write a Riemann sum approximating the total number of organisms in
        the petri dish.
      \end{problem}

      \begin{solution}
        Summing up the approximate number of organisms in each slice, the
        total number of organisms in the petri dish is approximately
        $\boxed{\sum_{k=1}^n f(x_k) \cdot 2 \pi x_k \Delta x}$. 
      \end{solution}

    \item
      \begin{problem}[\vfill]
        Write an integral that gives the exact number of organisms in the
        dish.
      \end{problem}

      \begin{solution}
        Taking the limit as the number of slices $\to \infty$, the exact
        number of organisms in the dish is $\boxed{\int_0^{9} f(x) \cdot
          2\pi x \,dx}$.
      \end{solution}

    \item
      \begin{problem}[\nstretch{3}]
        Find the number of organisms in the dish if $f(x)=100 e^{-x^2}$
        organisms per square centimeter.
      \end{problem}

      \begin{solution}
        In this case, the number of organisms is
        \begin{align*}
          \int_0^9 f(x) \cdot 2 \pi x\,dx &= 200 \pi \int_0^9 x
          e^{-x^2}\,dx \\
          \intertext{To evaluate this integral, we'll substitute $u =
            -x^2$.  Then, $du = -2x\,dx$, so $-\frac{du}{2} = x\,dx$:}
          &= 200 \pi \int_0^{-81} e^u \left( -\frac{du}{2} \right) \\
          &= -100 \pi \int_0^{-81} e^u\,du \\
          &= -100 \pi e^u \Big|_0^{-81} \\
          &= \boxed{-100 \pi e^{-81} + 100 \pi}
        \end{align*}
      \end{solution}

    \end{enumerate}

  \item Suppose that the density of organisms in another circular petri
    dish varies with the distance from a strip of nutrients running along
    the diameter of the dish. The density at a distance $x$ centimeters
    from the line of nutrients is given by $g(x)$ organisms per square
    centimeter. The petri dish is 18 centimeters in diameter.
    \begin{enumerate}

    \item
      \begin{problem}[\stretch{3}]
        Show in a sketch how you will slice the petri dish to solve this
        problem.
      \end{problem}

      \begin{solution}
        As usual, we want to slice the region so that the density in each
        slice is approximately constant.  In this case, density varies
        with distance to the nutrient strip, so we'll need to slice
        parallel to the strip.  If we draw the nutrient strip running
        vertically through the petri dish (it's shown in blue), then the
        situation looks like this:
        \begin{center}
          \begin{tikzpicture}
            \begin{axis}[x=1.5cm, y=1.5cm, xmin=0, xmax=1.2, xtick={0.01,
                1}, xticklabels={0, 9}, hide y axis, xlabel=$x$, ymin=-1,
              ymax=1, clip=false]
              \draw (0, 0) circle (1); %
              \pgfplotsforeachungrouped \x in {-0.875, -0.75, ..., 0.876}
              {
                \edef\temp{
                  \noexpand\draw (\x, {-sqrt(1-abs(\x)^2)}) -- (\x,
                  {sqrt(1-abs(\x)^2)}); %
                }
                \temp
              }
              \draw[blue, very thick] (0, -1) -- (0, 1);
            \end{axis}
          \end{tikzpicture}
        \end{center}
      \end{solution}

    \item
      \begin{problem}[\stretch{2}]
        Write a Riemann sum approximating the total number of organisms in
        the petri dish.
      \end{problem}

      \begin{solution}
        By symmetry, there will be the same amount of bacteria on either
        side of the nutrient strip.  In fact, symmetry tells us that, if
        we cut the petri dish both horizontally and vertically through the
        center, each quarter has the same amount of bacteria, so let's
        just focus on the top right quadrant of the dish.  Remembering
        that $x$ measures the distance from the nutrient strip, here is a
        picture of the quarter dish with the $x$-axis:
        \begin{center}
          \begin{tikzpicture}
            \begin{axis}[xmin=0, xmax=110, ymin=0, ymax=100, x=0.012in,
              y=0.012in, hide y axis, axis x line=bottom, clip=false,
              xlabel=$x$, xtick={0, 100}, xticklabels={0, 9}]  
              \begin{scope}[yshift=8pt]
                \draw (0, 0) -- (100, 0) arc[radius=100, start angle=0, end
                angle=90]; %
                \draw[blue, thick] (0, 0) -- (0, 100);
              \end{scope}
            \end{axis}
          \end{tikzpicture}
        \end{center}
        Slicing parallel to the nutrient strip means slicing the interval
        $[0, 9]$; as usual, we'll slice into $n$ pieces of equal
        thickness $\Delta x$ and label the values where we sliced as $x_0
        = 0, x_1, \ldots, x_n = 9$.  Let's focus on the $k$-th slice:
        \begin{center}
          \begin{tikzpicture}
            \begin{axis}[xmin=0, xmax=22, ymin=0, ymax=20, x=0.06in,
              y=0.06in, hide y axis, axis x line=bottom, clip=false,
              xlabel=$x$, xtick={0, 2.5, 5, ..., 20.1},
              xticklabels={$x_0$,,,,$x_k$,,,,$x_n$}] 
              \begin{scope}[yshift=8pt]
                \addplot+[domain=7.5:10, fill=cyan, draw=none,yshift=-8pt]
                {sqrt(20^2-\x^2)} -- (10, 0) -- (7.5, 0);
                \draw (0, 0) -- (20, 0) arc[radius=20, start angle=0, end
                angle=90]; %
                \draw[blue, thick] (0, 0) -- (0, 20);
                \pgfplotsforeachungrouped \x in {2.5, 5, ..., 19.9}
                {
                  \edef\temp{
                    \noexpand\draw (\x, 0) -- (\x, {sqrt(20^2-\x^2)});
                  }
                  \temp
                }              
              \end{scope}
            \end{axis}
          \end{tikzpicture}
        \end{center}
        The number of bacteria in the $k$-th slice is approximately
        $(\text{density in $k$-th slice})(\text{area of $k$-th slice})$.
        The density in the $k$-th slice is approximately $g(x_k)$.  To
        approximate the area of the $k$-th slice, we can approximate the
        slice as a rectangle of width $\Delta x$.  However, we still need
        to express its height in terms of $x_k$.  The key is to use the
        shape of the dish: it's part of a circle, so we need to use the
        radius.  If we draw in the radius (shown in red) below, then we
        can express the height $h_k$ in terms of $x_k$:
        \begin{center}
          \begin{tikzpicture}
            \begin{axis}[xmin=0, xmax=22, ymin=0, ymax=20, x=0.06in,
              y=0.06in, hide y axis, axis x line=bottom, clip=false,
              xlabel=$x$, xtick=\empty, hide x axis] 
              \begin{scope}[yshift=8pt]
                \addplot+[domain=7.5:10, fill=cyan, draw=none,yshift=-8pt]
                {sqrt(20^2-\x^2)} -- (10, 0) -- (7.5, 0); %
                \draw (7.5, 0) -- (7.5, {sqrt(20^2-7.5^2)}); %
                \draw (10, 0) -- node[right, inner sep=1pt] {$h_k$} (10,
                {sqrt(20^2-10^2)}); %
                \draw (0, 0) -- (20, 0) arc[radius=20, start angle=0, end
                angle=90]; %
                \draw[blue, thick] (0, 0) -- (0, 20); %
                \draw[decorate, decoration={brace, mirror}, yshift=-2pt] (0,
                0) -- node[yshift=-2pt, below] {$x_k$} (10, 0); %
                \draw[red, thick] (0, 0) -- (10, {sqrt(20^2-10^2)}); %
              \end{scope}
            \end{axis}
          \end{tikzpicture}
        \end{center}
        We now see a right triangle in the picture: the hypotenuse is the
        radius of the dish, which is 9; the two legs have lengths $x_k$
        and $h_k$.  Therefore, by the Pythagorean Theorem, $x_k^2 + h_k^2
        = 9^2$, so $h_k = \sqrt{81 - x_k^2}$.  So, we can approximate the
        area of this slice as $h_k \Delta x = \sqrt{81 - x_k^2} \Delta x$.
        Therefore,
        \[ \text{number of bacteria in $k$-th slice} \approx g(x_k) \cdot
        \sqrt{81 - x_k^2} \Delta x.\] Summing over all slices and
        remembering that we were only looking at a quarter of the dish,
        the number of bacteria in the dish is approximately $\boxed{4
          \sum_{k=1}^n g(x_k) \cdot \sqrt{81 - x_k^2} \Delta x}$.
      \end{solution}    

    \item
      \begin{problem}[\nstretch{2}]
        Write an integral that gives the exact number of organisms in the
        dish.
      \end{problem}

      \begin{solution}
        Taking the limit of our Riemann sum as the number of slices $\to
        \infty$, the number of bacteria in the dish is exactly
        $\boxed{4 \int_0^9 g(x) \cdot \sqrt{81 - x^2}\,dx}$. 
      \end{solution}
    \end{enumerate}

  \end{enumerate}

\item
\begin{grading}
    20 points. 4/8/8.
\end{grading}

  The Three Aces pizzeria has two chefs, Paul and Tim.  Both toss lots of
  garlic on their pizzas.
  \begin{itemize}
  \item Chef Paul tosses garlic so that the density of garlic on his pizza
    is given by
    \[ g(x) = \frac{1}{1 + x} \text{ grams per square inch of pizza},\]
    where $x$ is distance from the center of the pizza.
  \item Chef Tim tosses garlic so that the density of garlic on his pizza
    is given by
    \[ g(x) = \frac{1}{1 + x} \text{ grams per square inch of pizza},\]
    where $x$ is distance from the \emph{edge} of the pizza.
  \end{itemize}
  Each pizza is 14 inches in diameter.

  \begin{enumerate}
  \item
    \begin{problem}[\vfill]
      Which chef puts more garlic on a pizza?  (You should be able to
      figure this out without doing any calculations; think about which
      parts of each chef's pizza are most densely covered with garlic and
      which parts are least densely covered.)  Explain.
    \end{problem}

    \begin{solution}
      Both pizzas are the same size and have density of garlic varying
      from $\frac{1}{8}$ grams of garlic per square inch to $1$ gram of
      garlic per square inch.  The difference is that, on Chef Tim's
      pizza, the edges of the pizza are where the garlic is most dense
      while, on Chef Paul's pizza, the center of the pizza is where the
      garlic is most dense.  Since there is a lot more edge than center,
      \fbox{Chef Tim}'s pizza ends up with more garlic.  Here's a
      visualization, where the dots represent garlic:
      \begin{center}
        \begin{tabular}{c c c}
          \includegraphics[width=1.25in]{images/pizza-garlic-variations-1.pdf}
          & \hspace{0.25in} &
          \includegraphics[width=1.25in]{images/pizza-garlic-variations-2.pdf}
          \\
          Chef Paul's pizza & & Chef Tim's pizza
        \end{tabular}
      \end{center}
    \end{solution}

  \item
    \begin{problem}[\nstretch{2}]
      Find the amount of garlic on each of Chef Paul's pizzas.
    \end{problem}

    \begin{solution}
      We must slice the pizza so that the density of garlic within each
      slice is approximately constant.  Since density varies with distance
      to the center, we slice into concentric rings.  So, we slice the
      interval $[0, 7]$ (in $x$) into $n$ pieces of equal width $\Delta x
      = \frac{7}{n}$, label $x_0 = 0, \ldots, x_n = 7$ as usual (so $x_k =
      k \Delta x$), and focus on the $k$-th slice:
      \begin{center}
          \begin{tikzpicture}
            \begin{axis}[x=0.007in, y=0.007in, xmin=0, xmax=115,
              ymin=-100, ymax=100, axis on top=true, clip=false, hide y
              axis, xtick={0.1, 75, 100}, xticklabels={$x_0$, $x_k$, $x_n$},
              xlabel=$x$, tick label style={fill=none}]              
              \draw[fill=cyan, draw=none] (0, 0) circle(75); %
              \draw[fill=white, draw=none] (0, 0) circle (62.5); %
              \pgfplotsforeachungrouped \x in {12.5, 25, ..., 100.1}
              {
                \edef\temp{
                  \noexpand\draw (0, 0) circle(\x);
                }
                \temp
              }              
            \end{axis}
          \end{tikzpicture}
        \end{center}
      The area of this ring is approximately $2 \pi x_k \Delta x$ cm$^2$,
      and the density of garlic in this ring is approximately $g(x_k)$
      g/in$^2$, so the amount of garlic in this ring is approximately $2
      \pi x_k g(x_k) \Delta x$ g.  Summing over all slices and taking a
      limit as the number of slices $\to \infty$, the amount of garlic in
      the pizza is $\displaystyle \int_0^7 2 \pi x g(x)\,dx$.  Now, we
      must evaluate the integral:
      \begin{align*}
        \int_0^7 2 \pi x g(x)\,dx &= 2 \pi \int_0^7 \frac{x}{1 + x}\,dx
        \\
        \intertext{Let's substitute $u = 1 + x$; then $du = dx$ and $x = u
          - 1$.  Also, $u$ goes from $1 + 0 = 1$ to $1 + 7 = 8$:}
        &= 2 \pi \int_1^{8} \frac{u - 1}{u}\,du \\
        &= 2 \pi \int_1^8 \left( 1 - \frac{1}{u} \right) du \\
        &= 2 \pi \left( u - \ln |u| \right) \Big|_1^8 \\
        &= \boxed{2 \pi (7 - \ln 8) \text{ grams}}
      \end{align*}
      (This is about $31$ grams.)
    \end{solution}

  \item
    \begin{problem}[\vfill]
      Find the amount of garlic on each of Chef Tim's pizzas.      
    \end{problem}

    \begin{solution}
      Again, we must slice the pizza so that the density of garlic within
      each slice is approximately constant.  Since density varies with
      distance to the edge, we still slice into concentric rings.  So
      what's different?  Well, the meaning of $x$ is now different: $x$ is
      0 on the edge of the pizza, and $x$ is 7 at the center of the pizza.

     If we let $r$ represent the distance from the center of the pizza and we sliced into pieces of width $\Delta r$,
     like this:
      \begin{center}
        \begin{tikzpicture}
          \begin{axis}[x=0.007in, y=0.007in, xmin=0, xmax=115,
            ymin=-100, ymax=100, axis on top=true, clip=false, hide y
            axis, xtick={0.1, 75, 100}, xticklabels={$r_0$, $r_k$, $r_n$},
            xlabel=$r$, tick label style={fill=none}]              
            \draw[fill=cyan, draw=none] (0, 0) circle(75); %
            \draw[fill=white, draw=none] (0, 0) circle (62.5); %
            \pgfplotsforeachungrouped \x in {12.5, 25, ..., 100.1}
            {
              \edef\temp{
                \noexpand\draw (0, 0) circle(\x);
              }
              \temp
            }              
          \end{axis}
        \end{tikzpicture}
      \end{center}
      then the area of the $k$-th concentric ring would be $2\pi r_k\Delta r$,
      same as before. However, since our density is in terms of $x$ and we want 
      our
      integral to be in terms of $x$, we should write the area in terms of $x$.

      What is the relationship between $x$ and $r$? Since $x$ is the distance
      from the edge and $r$ is the distance from the center, $r=7-x$.
      This means that, if we slice into pieces of width $\Delta x$, the area of the $k$-th concentric ring
      would be $2\pi(7- x_k)\Delta x$. 
      
      Summing over all slices and taking a limit as the number
      of slices $\to \infty$, the amount of garlic in the pizza is
      $\displaystyle \int_0^7 2 \pi (7-x) g(x)\,dx$.  Now, we must
      evaluate the integral:
      \begin{align*}
        \int_0^7 2 \pi (7-x) g(x)\,dx &= \int_0^7 \frac{2 \pi (7-x)}{1 + x}\,dx \\
        &= 2 \pi \int_0^7 \frac{7-x}{1+x}\,dx \\
        \intertext{Let's substitute $u = 1+x$; then $du = dx$ and $x =
          u-1$, and $u$ goes from $1 + 0 = 1$ to $1 + 7 = 8$:} %
        &= 2 \pi \int_1^8 \frac{7- (u-1)}{u} du \\
        &= 2 \pi \int_1^8 \frac{8-u}{u} du \\
        &= 2 \pi \int_1^8 \left(\frac{8}{u} - 1 \right) du \\
        &= 2 \pi \Bigl[8 \ln |u| - u \Bigr]_1^8 \\
        &= \boxed{2 \pi \left( 8 \ln 8 - 7 \right)}
      \end{align*}
      (This is about $60.5$ grams, which confirms our intuition that
      Chef Tim's pizza has more garlic.)
    \end{solution}

  \end{enumerate}

% \item ({Problem Set 31}, {{\#3}})

%   \begin{minipage}[t]{2in}
%     \raisebox{2ex-\height}{      
%       \includegraphics[width=1.95in]{images/bompas-parr-jelly-cone}}
%     \begin{center}
%       \parbox[t]{1.75in}{\emph{\small This cone is missing its tip, but Sam
%           and Harry's cones do have tips!}}
%     \end{center}
%   \end{minipage}
%   \begin{minipage}[t]{\linewidth-2.1in}
%     \setlength{\parskip}{8pt} Bompas \& Parr is a British company that
%     specializes in artwork made from gelatin.  One recipe they have
%     published is for a cone-shaped ``striped jelly''.

%     Sam and Harry each decide to make a jelly cone with a radius of 4 cm
%     and height 7 cm; the tip of each cone will point up.  The density of
%     sugar varies in the finished cones.  In Harry's cone, the density of
%     sugar is given by $f(x) = x$ mg/cm$^3$, where $x$ is the vertical
%     distance (in cm) from the tip of the cone.  In Sam's cone, the density
%     of sugar is given by $f(x) = x$ mg/cm$^3$, where $x$ is the vertical
%     distance (in cm) from the \emph{bottom} of the cone.
%   \end{minipage}

%   \begin{enumerate}
%   \item
%     \begin{problem}
%       Whose cone will have more sugar in it?  (You should be able to
%       figure this out without doing any calculations.)
%     \end{problem}

%     \begin{solution}
%       The two cones are the same size, and both have density of sugar
%       varying from $0$ mg/cm$^3$ to $7$ mg/cm$^3$.  The difference is
%       that, in Harry's cone, the bottom is where the sugar is more dense; in
%       Sam's cone, the top is where the sugar is more dense.  Since there
%       is more bottom than top, \fbox{Harry's} cone will have more sugar in
%       it.
%     \end{solution}

%   \item
%     \begin{problem}
%       Find the amount of sugar in Harry's cone.
%     \end{problem}

%     \begin{solution}
%       Since the density of sugar varies with vertical distance to the tip,
%       we should slice the cone parallel to its base so that the density of
%       sugar within each slice is approximately constant.  Since $x =
%       0$ at the top of the cone and $x = 7$ at the bottom, we're slicing
%       the interval $[0, 7]$; we'll label $x_0 = 0, \ldots, x_n = 7$ as
%       usual. 
%       \begin{center}
%         \begin{tikzpicture}[x=0.5cm, y=0.5cm]
%           \coordinate (right) at (4, 0);
%           \coordinate (left) at (-4, 0);
%           \coordinate (top) at (0, 7);

%           \begin{axis}[xmin=-4, xmax=4.5, ymin=-1, ymax=7, clip=false, hide x
%             axis, axis y line=right, x=0.5cm, y=0.5cm, ytick={0, 0.7, ...,
%               7.01},
%             yticklabels={$x_n = 7$,,,,$x_k$,,,,,,$x_0 = 0$}, 
%             axis line style={<-}]
%             \draw[fill=cyan, draw=none] (2, 3.5) arc[x radius=2,
%             y radius=0.5, start angle=0, end angle=-180] -- (-2.4, 2.8)
%             arc[x radius=2.4, y radius=0.6, start angle=180, end
%             angle=360] -- cycle; %
%             \draw (right) arc[x radius=4, y radius=1, start angle=0, end
%             angle=-180]; %
%             \draw (right) -- (top) -- (left); %
%             \draw (4.5, -1) node[below] {$x$}; %
%             \pgfplotsforeachungrouped \y in {0, 0.7, ..., 7.01} {
%               \pgfmathsetmacro{\x}{4/7*(7-\y)} %
%               \edef\temp{%
%                 \noexpand\draw (\x, \y) arc[x radius=\x, y radius={\x/4},
%                 start angle=0, end angle=-180];
%               }              
%               \temp
%             }
%           \end{axis}
%         \end{tikzpicture}
%       \end{center}
%       Then,
%       \[ \text{amount of sugar in $k$-th slice} \approx (\text{density of
%         sugar in $k$-th slice})(\text{volume of $k$-th slice}).\]
%       The density of sugar in the $k$-th slice is approximately $f(x_k)$.
%       To find the volume of the $k$-th slice, we can approximate the slice
%       as a disk; then, its volume is approximately $\pi (\text{radius})^2
%       \cdot \Delta x$, but we need to express its radius in terms of
%       $\Delta x$.  Here's a picture that helps us do that: 
%       \begin{center}
%         \begin{tikzpicture}[x=0.5cm, y=0.5cm]
%           \coordinate (right) at (4, 0);
%           \coordinate (left) at (-4, 0);
%           \coordinate (top) at (0, 7);

%           \begin{axis}[xmin=-4, xmax=4.5, ymin=-1, ymax=7, clip=false, hide x
%             axis, axis y line=right, x=0.5cm, y=0.5cm, ytick={0, 0.7, ...,
%               7.01},
%             yticklabels={$x_n = 7$,,,,$x_k$,,,,,,$x_0 = 0$}, 
%             axis line style={<-}]
%             \draw[fill=cyan, draw=none] (2, 3.5) arc[x radius=2,
%             y radius=0.5, start angle=0, end angle=-180] -- (-2.4, 2.8)
%             arc[x radius=2.4, y radius=0.6, start angle=180, end
%             angle=360] -- cycle; %
%             \draw (right) arc[x radius=4, y radius=1, start angle=0, end
%             angle=-180]; %
%             \draw (right) -- (top) -- (left); %
%             \draw (4.5, -1) node[below] {$x$}; %
%             \pgfplotsforeachungrouped \y in {2.8, 3.5} {
%               \pgfmathsetmacro{\x}{4/7*(7-\y)} %
%               \edef\temp{%
%                 \noexpand\draw (\x, \y) arc[x radius=\x, y radius={\x/4},
%                 start angle=0, end angle=-180];
%               }              
%               \temp
%             }%            
%             \draw[dashed] (2, 3.5) arc[x radius=2, y radius=0.5,
%             start angle=0, end angle=180]; %
%             \draw[red, thick] (0, 7) -- (0, 0) -- (4, 0);
%             \draw[red, thick] (0, 2.8) -- (2.4, 2.8);
%           \end{axis}
%         \end{tikzpicture}
%       \end{center}
%       In the picture, we see similar triangles; the larger triangle has
%       height 7 and width 4, while the smaller triangle has height $x_k$,
%       and its width is exactly the radius we're interested in.  So,
%       $\displaystyle \frac{7}{4} = \frac{x_k}{\text{approximate radius of
%           $k$-th slice}}$.  Therefore, the approximate radius of the
%       $k$-th slice is $\frac{4 x_k}{7}$.  Putting this all together, the
%       amount of sugar in the $k$-th slice is approximately $\displaystyle
%       f(x_k) \cdot \pi \left( \frac{4 x_k}{7} \right)^2 \Delta x \approx
%       \frac{16\pi}{49}x_k^2\Delta x$.  Adding these up and taking the
%       limit as the number of slices $\to \infty$, the amount of sugar in
%       the entire cone is
%       \begin{align*}
%         \boxed{\int_0^7 \frac{16\pi}{49}x^2 f(x)\,dx} &=
%         \frac{16\pi}{49} \int_0^7 x^3\,dx \\
%         &=
%         \left. \frac{16\pi}{49}\cdot\frac{1}{4} x^4 \right|_0^7 \\
%         &= \boxed{196\pi \text{ mg of sugar}}
%       \end{align*}
%     \end{solution}

%   \item
%     \begin{problem}
%       Find the amount of sugar in Sam's cone. 
%     \end{problem}

%     \begin{solution}
%       Since the density of sugar varies with vertical distance from the
%       bottom of the cone, we should again slice parallel to the base of
%       the cone so that the density of sugar within each slice is
%       approximately constant.  Now, $x = 0$ at the bottom of the cone, and
%       $x = 7$ at the top, so we're again slicing the interval $[0, 7]$;
%       however, the $x$-axis is oriented differently:
%       \begin{center}
%         \begin{tikzpicture}[x=0.5cm, y=0.5cm]
%           \coordinate (right) at (4, 0);
%           \coordinate (left) at (-4, 0);
%           \coordinate (top) at (0, 7);

%           \begin{axis}[xmin=-4, xmax=4.5, ymin=-0.02, ymax=8, clip=false,
%             hide x axis, axis y line=right, x=0.5cm, y=0.5cm, ytick={7,
%               6.3, ..., -0.01}, yticklabels={$x_n = 7$,,,,,,$x_k$,,,,$x_0
%               = 0$}]
%             \draw[fill=cyan, draw=none] (2, 3.5) arc[x radius=2, y
%             radius=0.5, start angle=0, end angle=-180] -- (-2.4, 2.8)
%             arc[x radius=2.4, y radius=0.6, start angle=180, end
%             angle=360] -- cycle; %
%             \draw (right) arc[x radius=4, y radius=1, start angle=0, end
%             angle=-180]; %
%             \draw (right) -- (top) -- (left); %
%             \draw (4.5, 8) node[above] {$x$}; %
%             \pgfplotsforeachungrouped \y in {0, 0.7, ..., 7.01} {
%               \pgfmathsetmacro{\x}{4/7*(7-\y)} %
%               \edef\temp{%
%                 \noexpand\draw (\x, \y) arc[x radius=\x, y radius={\x/4},
%                 start angle=0, end angle=-180]; } \temp }
%           \end{axis}
%         \end{tikzpicture}
%       \end{center}
%       Again,
%       \[ \text{amount of sugar in $k$-th slice} \approx (\text{density of
%         sugar in $k$-th slice})(\text{volume of $k$-th slice}),\] and the
%       density of sugar in the $k$-th slice is approximately $f(x_k)$.  The
%       $k$-th slice is approximately a disk, and we can again use similar
%       triangles to approximate its radius:
%       \begin{center}
%         \begin{tikzpicture}[x=0.5cm, y=0.5cm]
%           \coordinate (right) at (4, 0);
%           \coordinate (left) at (-4, 0);
%           \coordinate (top) at (0, 7);

%           \begin{axis}[xmin=-4, xmax=4.5, ymin=-0.01, ymax=8, clip=false,
%             hide x axis, axis y line=right, x=0.5cm, y=0.5cm, ytick={7,
%               6.3, ..., -0.001}, yticklabels={$x_n = 7$,,,,,,$x_k$,,,,$x_0
%               = 0$}]
%             \draw[fill=cyan, draw=none] (2, 3.5) arc[x radius=2,
%             y radius=0.5, start angle=0, end angle=-180] -- (-2.4, 2.8)
%             arc[x radius=2.4, y radius=0.6, start angle=180, end
%             angle=360] -- cycle; %
%             \draw (right) arc[x radius=4, y radius=1, start angle=0, end
%             angle=-180]; %
%             \draw (right) -- (top) -- (left); %
%             \draw (4.5, 8) node[above] {$x$}; %
%             \pgfplotsforeachungrouped \y in {2.8, 3.5} {
%               \pgfmathsetmacro{\x}{4/7*(7-\y)} %
%               \edef\temp{%
%                 \noexpand\draw (\x, \y) arc[x radius=\x, y radius={\x/4},
%                 start angle=0, end angle=-180];
%               }              
%               \temp
%             }%            
%             \draw[dashed] (2, 3.5) arc[x radius=2, y radius=0.5,
%             start angle=0, end angle=180]; %
%             \draw[red, thick] (0, 7) -- (0, 0) -- (4, 0);
%             \draw[red, thick] (0, 2.8) -- (2.4, 2.8);
%           \end{axis}
%         \end{tikzpicture}
%       \end{center}
%       The larger triangle still has height 7 and width 4, while the
%       smaller triangle has height $7 - x_k$, and its width is exactly the
%       radius we're interested in.  So, $\displaystyle \frac{7}{4} =
%       \frac{7 - x_k}{\text{approximate radius of $k$-th slice}}$.
%       Therefore, the approximate radius of the $k$-th slice is $\frac{4 (7
%         - x_k)}{7}$, so the amount of sugar in the $k$-th slice is
%       approximately $\displaystyle f(x_k) \cdot \pi \left( \frac{4 (7 -
%           x_k)}{7} \right)^2 \Delta x \approx \frac{16 \pi}{49} (7 - x_k)^2
%       \Delta x$.  Summing over all slices and taking the limit as the
%       number of slices $\to \infty$, the amount of sugar in the entire
%       cone is
%       \begin{align*}
%         \boxed{\int_0^7 \frac{16\pi}{49} (7 - x)^2 f(x)\,dx} &=
%         \frac{16\pi}{49} \int_0^7 (7 - x)^2 x\,dx \\
%         &= \frac{16\pi}{49} \int_0^7 (49 x - 14 x^2 + x^3)\,dx \\
%         &= \frac{16 \pi}{49} \left( \left. \frac{49}{2} x^2 - \frac{14}{3}
%             x^3 + \frac{1}{4} x^4 \right|_0^7 \right) \\
%         &= \boxed{\frac{196 \pi}{3} \text{ mg of sugar}}
%       \end{align*}
%       This confirms our intuition that Harry's cone has more sugar. 
%     \end{solution}

%   \end{enumerate}

% \item

% \begin{problem}[\nstretch{1.5}]
%     A boy left a puddle of maple syrup on a picnic table after a
%     pancake breakfast.  The maple syrup is lying in a circular puddle of
%     diameter 3 inches (radius 1.5 inches). Five hours after breakfast, the
%     density of ants varies with the distance from the edge of the puddle
%     of syrup. (There are no ants {\it in} the syrup, only around it.)  The
%     density of ants is given by $\rho(x)$ ants per square inch, where $x$
%     is the distance from the edge of the puddle. Write an integral for the
%     number of ants within 7 inches of the edge of the puddle.
%   \end{problem}

%   \begin{solution}
%     Here is a sketch of the region we are interested in: 
%     \begin{center}
%       \begin{tikzpicture}
%         \begin{axis}[xmin=3, xmax=10, ymin=-9, ymax=9, clip=false, x=0.2cm,
%           y=0.2cm, hide y axis, xtick={3.01, 9}, 
%           xticklabels={0, 7}]
%           \draw[fill=lightgray, fill opacity=0.5] (0, 0) circle (9);
%           \draw[fill=white] (0, 0) circle (3);
%           \filldraw[black] (0, 0) circle(1.5pt);
%           \draw (10, 0) node[right] {$x$};
%           \draw (0, 0) -- node[anchor=210, inner sep=0.5pt] {1.5} (120: 3);
%           % \draw (0, 0) -- node[Anchor=225, inner sep=1pt] {9} (135:9);
%         \end{axis}
%       \end{tikzpicture}
%     \end{center}
%     The shaded portion is the area we need to look at; the hole in the
%     middle is the puddle of maple syrup.  We've also drawn an ``$x$-axis''
%     in the picture to show that $x$ measures distance from the edge of the
%     puddle.  In particular, $x = 0$ at the edge of the puddle, and $x = 7$
%     at the outer edge of the area we're interested in. 

%     Since the density of ants varies with distance from the edge of the
%     puddle, we need to slice so that all points in a slice are (roughly)
%     the same distance from the edge of the puddle.  So, we'll slice into
%     concentric rings around the puddle.  This means we'll slice $[0, 7]$
%     into $n$ pieces of equal length $\Delta x = \frac{7}{n}$, label $x_0 =
%     0, \ldots, x_n = 7$ as usual (so $x_k = k \Delta x$), and focus on the
%     $k$-th slice:
%     \begin{center}
%       \begin{tikzpicture}
%         \begin{axis}[xmin=3, xmax=10, ymin=-9, ymax=9, clip=false, x=0.2cm,
%           y=0.2cm, hide y axis, xtick={3, 4, ..., 9}, 
%           xticklabels={,,,,$x_k$,,}, axis on top=true, tick label
%           style={fill=none}] 
%           \draw[fill=white!90!black] (0, 0) circle (9);
%           \draw[fill=cyan] (0, 0) circle (7);
%           \draw[fill=white!90!black] (0, 0) circle (6);
%           \draw[fill=white] (0, 0) circle (3);
%           \filldraw[black] (0, 0) circle(1.5pt);
%           \draw (10, 0) node[right] {$x$};
%           \draw (0, 0) -- node[anchor=210, inner sep=0.5pt] {1.5} (120: 3);
%           \pgfplotsforeachungrouped \x in {3, 4, ..., 9}
%           {
%             \edef\temp{
%               \noexpand\draw (0, 0) circle(\x);
%             }
%             \temp
%           }              
%         \end{axis}
%       \end{tikzpicture}
%     \end{center}
%     This slice has thickness $\Delta x$ and outer radius $r_k = 1.5+x_k$.
%     Therefore, we can approximate its area as $2 \pi r_k \Delta x = 2 \pi
%     (1.5 + x_k) \Delta x$.  The density of ants in this slice is
%     approximately $\rho(x_k)$, so the number of ants in this slice is
%     approximately $\rho(x_k) \cdot 2 \pi (1.5 + x_k) \Delta x$.  Summing
%     over all slices and taking the limit as the number of slices $\to
%     \infty$, the number of ants within $7$ inches of the edge of the
%     puddle is $\boxed{\int_0^{7} 2\pi(1.5+x) \rho(x)\,dx}$.
%   \end{solution}  

\item 
\begin{problem}[\vfill\newpage]
Pizza Pythagoras is known for its ``pizza wedges,'' which are shaped
like right triangles with sides of 3, 4, and 5 inches.  Parmesan
cheese is sprinkled on each wedge so that the density of cheese is
given by $\rho(x)$ ounces per square inch, where $x$ measures distance
(in inches) from the 3-inch side of the slice. 
Write a sum
(involving $n$ slices) that approximates the amount of Parmesan on
each wedge. Hint: consider graphing the wedge on the $xy$-plane first before slicing.

Be sure to explain clearly how you are slicing and how you obtain your
approximation for the amount of cheese in the $k$-th slice.  If you
use notation such as $\Delta x$ or $x_k$, be sure to explain exactly
what it means.
\end{problem}

\begin{solution}
Let's imagine the slice as sitting on the coordinate plane so that the
3-inch side lies along the $y$-axis and the 4-inch side lies along the
$x$-axis.
\begin{center}
\includegraphics[width=1.5in]{images/pizza-pythagoras-region}
\end{center}
As in the previous problem, this choice has the advantage that $x$ in
the picture matches $x$ in the problem: it measures distance from the
3-inch side of the slice.  Notice that the 5-inch side of the pizza
connects the points $(0, 3)$ and $(4, 0)$; therefore, it can be
described by the line $y = - \frac{3}{4} x + 3$.

Now, we must decide how to slice the pizza wedge.  Since we
\textbf{want the density of Parmesan within each slice to be
approximately constant} and we know that the density varies with
distance to the 3-inch side, we should slice vertically.  Let's slice
the wedge into $n$ vertical strips, each of width $\Delta x =
\frac{4}{n}$.  We'll label the resulting endpoints $x_0 = 0, x_1 =
\Delta x, \ldots, x_n = 4$.  Then, a general formula for $x_k$ is $x_k
= k \Delta x$.
\begin{center}
\includegraphics[width=1.5in]{images/pizza-pythagoras-slice}
\end{center}
The density of Parmesan in the $k$-th slice is approximately
$\rho(x_k)$.  (Really, $\rho(x_k)$ is the density of Parmesan along
the right edge of the slice, but we've sliced so that the density
within a slice doesn't vary too much, so this is a reasonable
approximation.)  We can approximate the area of the $k$-th slice by
drawing a rectangle that approximates the slice:
\begin{center}
\includegraphics[width=1.5in]{images/pizza-pythagoras-approximating-rectangle}
\end{center}
This rectangle has width $\Delta x$ and height $- \frac{3}{4} x_k +
3$, so its area is $\left( - \frac{3}{4} x_k + 3 \right) \Delta x$.
To summarize, so far we've made the following approximations:
\begin{align*}
\textrm{Parmesan density in $k$-th slice} & \approx \rho(x_k) \\
\textrm{area of $k$-th slice} & \approx \left( - \frac{3}{4} x_k + 3
\right) \Delta x \\
\shortintertext{So,}
\textrm{amount of Parmesan in $k$-th slice} & \approx
(\textrm{Parmesan density in $k$-th slice}) \cdot (\textrm{area of
$k$-th slice}) \\
& \approx \rho(x_k) \cdot \left( - \frac{3}{4} x_k + 3 \right)
\Delta x
\end{align*}
To get an approximation for the amount of Parmesan on the whole wedge
of pizza, we must add up the approximate amounts of Parmesan on each
slice, which gives us the approximation \[\boxed{\rho(x_1) \cdot
\left( - \frac{3}{4} x_1 + 3 \right) \Delta x + \rho(x_2) \cdot
\left( - \frac{3}{4} x_2 + 3 \right) \Delta x + \cdots + \rho(x_n)
\cdot \left( - \frac{3}{4} x_n + 3 \right) \Delta x}\] or
$\boxed{\sum_{k=1}^n \rho(x_k) \cdot \left( - \frac{3}{4} x_k + 3
\right) \Delta x}$.
\end{solution}

\item

\begin{grading}
    1 point each.
\end{grading}

 \textbf{Reflect Back.} \emph{Which way to
    slice?}
Below are two problems about cylinders.  In each part, decide whether
  you should slice into disks or tubes to solve the problem.  (You do
  \textbf{not} need to actually solve the problems!)
  \begin{center}
    \begin{tabular}{c c c}
      \includegraphics[width=1.25in]{images/density-slicing-cylinder-ways-disks.pdf}
      & 
      \hspace{0.25in}
      & \includegraphics[width=1.25in]{images/density-slicing-cylinder-ways-tubes.pdf} 
      \\
      Disks & & Tubes 
    \end{tabular}
  \end{center}

  \begin{enumerate}
  \item
    \begin{problem}[\vfill\newpage]
      A chocolate popsicle is being made by taking a popsicle stick and
      dipping it repeatedly into chocolate and other tasty ingredients.
      Between each dip the ingredients are allowed to cool. What is
      constructed at the end of this process is a cylindrical popsicle
      with height 30 mm and radius 15 mm. Due to the way it has been made,
      the number of calories per cubic millimeter varies with $x$, where
      $x$ is the distance (in millimeters) from the central axis of the
      popsicle. Suppose that the density of calories in the popsicle is
      given by $\rho(x)$ calories per mm$^3$. Write an integral that gives
      the total number of calories in the popsicle.  (You may assume that
      the popsicle stick's width is negligible.)
    \end{problem}

    \begin{solution}
      We need to slice into \fbox{tubes} in order for the density within
      any slice to be approximately constant. 
    \end{solution}

  \item
    \begin{problem}[\vfill]
      A scientist studying sea ice is analyzing a cylindrical sample of
      ice, called a ``core.'' It has height 2 meters and diameter 0.1
      meters.  Suppose that algae is distributed throughout the ice with
      density given by $\rho(h) = 300 + 50h^2$ g/m$^3$ where $h$ is the
      depth (in meters) from the top surface of the ice sample.  Find the
      number of grams of algae in the sample.
    \end{problem}

    \begin{solution}
      We need to slice into \fbox{disks} in order for the density within
      any slice to be approximately constant. 
    \end{solution}

  \end{enumerate}

\end{enumerate}

\spaceonly{\psreflection{
Optional reading for next class is \S 15.2, 31.2 in Gottlieb.
}}
\end{document}